\documentclass{article}
\usepackage{amsmath, amssymb, amsthm}
\usepackage{hyperref}
\usepackage{geometry}
\geometry{margin=1in}

\title{Erd\H{o}s Problem \#62}
\date{Last updated: 2025-08-31}
\author{Source: erdosproblems.com}

\begin{document}
\maketitle

\noindent\textbf{Status:} OPEN \quad \textbf{No prize}

\noindent\textbf{Tags:} graph theory

\medskip
\noindent\textbf{Problem Statement:}

If $G_1,G_2$ are two graphs with chromatic number $\aleph_1$ then must there exist a graph $G$ whose chromatic number is $4$ (or even $\aleph_0$) which is a subgraph of both $G_1$ and $G_2$?

\bigskip
\noindent\textbf{Remarks:}

Erd\H{o}s also asked \cite{Er87} about finding a common subgraph $H$ (with chromatic number either $4$ or $\aleph_0$) in any finite collection of graphs with chromatic number $\aleph_1$. 

Every graph with chromatic number $\aleph_1$ contains all sufficiently large odd cycles (which have chromatic number $3$), see [594]. This was proved by Erd\H{o}s, Hajnal, and Shelah \cite{EHS74}. Erd\H{o}s wrote \cite{Er87} that 'probably' every graph with chromatic number $\aleph_1$ contains as subgraphs all graphs with chromatic number $4$ with sufficiently large girth.

\bigskip
\noindent\textbf{References:}

[EHS74] Erd\H{o}s, P. and Hajnal, A. and Shelah, S., On some general properties of chromatic numbers. Topics in topology (Proc. Colloq., Keszthely, 1972) (1974), 243-255.
 
 [Er87] Erd\H{o}s, P., Some problems on finite and infinite graphs. Logic and combinatorics (Arcata, Calif., 1985) (1987), 223-228.

\bigskip
\noindent\small{Source: \url{https://www.erdosproblems.com/62}}
\end{document}
